\section{复变积分与变换的可计算桥梁}\label{deep:integral_complex_deep}

\begin{examtip}本节先建立公式，再展示实积分、卷积与微分方程的联系。\end{examtip}

\subsection{解析性、Cauchy公式与留数}\label{deep:integral_complex_deep:complex}

\begin{method}{解析性、Cauchy公式与留数}\textbf{适用条件} 函数在围道及内部解析，或只有有限个孤立极点。\par\textbf{步骤与依据} 辨认区域、方向和奇点阶数；用局部Laurent系数计算留数；闭合实轴围道时估计大弧。\par\textbf{易错点} 解析性不能由连续性替代；路径穿过奇点时原定理不适用。\end{method}

\begin{definition}{复导数与Laurent系数}复导数要求商$\frac{f(z+h)-f(z)}{h}$在复方向$h\to0$有同一极限。解析函数局部有Taylor展开；环域解析函数有Laurent展开$\sum_{n\in\mathbb Z}a_n(z-z_0)^n$，$a_{-1}$称留数。若$f=\frac{g}{(z-z_0)^m}$且$g$解析，则留数$\frac{g^{(m-1)}(z_0)}{(m-1)!}$。\end{definition}

\begin{theorem}{Cauchy积分公式及导数式}正向简单闭曲线$\Gamma$及内部含于解析区域，$z_0$在内部，则$f^{(n)}(z_0)=\frac{n!}{2\pi i}\int_\Gamma \frac{f(z)}{(z-z_0)^{n+1}}\dd z$。\end{theorem}\begin{proof}先对$n=0$，挖去$z_0$的小圆，Cauchy定理把外边积分化为小圆积分。写$f(z)=f(z_0)+[f(z)-f(z_0)]$，常数部分积分$2\pi if(z_0)$，余项由连续性趋0。对$z_0$在紧子区域内逐次求导，分母与围道有正距离，故可在积分号内求导，得到$n$阶式。\end{proof}

\begin{example}{竞赛拓展｜高阶极点与Cauchy公式}{难度：$\star\star\star$}\label{deep:ex:integral_complex_deep:cauchy}\textbf{题目} 计算正向圆$|z|=2$上的$\int\frac{\e^z}{(z-1)^3}\dd z$。\par\textbf{思路} 分母三次对应二阶导数。\par\textbf{详细解答} 取$f(z)=\e^z,z_0=1,n=2$，积分$=\frac{2\pi i f''(1)}{2!}=\pi i\e$。第二法在$z=1$展开$\e^z=\e[1+(z-1)+\frac{(z-1)^2}{2}+\cdots]$，留数为$\frac{\e}{2}$，同得$\pi i\e$。\par\textbf{方法点评} 导数公式最快，Laurent法有利于多个相乘因子。\par\textbf{变式与易错点} 分母幂$m$对应$m-1$阶导数。\end{example}

\begin{example}{竞赛拓展｜四次有理积分的留数法}{难度：$\star\star\star$}\label{deep:ex:integral_complex_deep:residue4}\textbf{题目} 计算$\int_{-\infty}^{\infty}\frac{\dd x}{1+x^4}$。\par\textbf{思路} 上半平面两个极点，弧长乘最大值控制大弧。\par\textbf{详细解答} 极点$z_1=\e^{\frac{i\pi}{4}},z_2=\e^{\frac{3i\pi}{4}}$，留数$\frac{1}{4z_j^3}$，和$=-\frac{i}{2\sqrt2}$。半径$R$大弧积分绝对值$\le\frac{\pi R}{R^4-1}\to0$。留数定理给实积分$2\pi i[-\frac{i}{2\sqrt2}]=\frac{\pi}{\sqrt2}$。由偶性半轴积分$\frac{\pi}{2\sqrt2}$。也可用前节四次分母的初等原函数在无穷两端取极限核对。\par\textbf{方法点评} 留数法免去实部分分式的长计算，初等法验证实值和符号。\par\textbf{变式与易错点} 必须取上半平面全部两个极点。\end{example}

\subsection{Fourier与Laplace的归一化、卷积和初值}\label{deep:integral_complex_deep:transforms}

\begin{method}{Fourier与Laplace的归一化、卷积和初值}\textbf{适用条件} 先选变换约定并验证绝对可积或指数增长控制。\par\textbf{步骤与依据} 全书实轴Fourier统一采用$\widehat f(\xi)=\int_\R f(x)\e^{-2\pi i\xi x}\dd x$；在反演成立的函数类上，$f(x)=\int_\R\widehat f(\xi)\e^{2\pi i\xi x}\dd\xi$，无额外系数。Laplace采用$\mathcal Lf(s)=\int_0^\infty\e^{-st}f(t)\dd t$。\par\textbf{易错点} 公式中的$2\pi$与初值项由约定决定。\end{method}

\begin{example}{竞赛拓展｜指数核、余弦变换与卷积}{难度：$\star\star\star$}\label{deep:ex:integral_complex_deep:fourierexp}\textbf{题目} 计算$f(x)=\e^{-a|x|}$的Fourier变换，$a>0$，并求$f*f$。\par\textbf{思路} 偶性化为余弦积分；卷积按$t<0,0<t<x,t>x$分段。\par\textbf{详细解答} $\widehat f(\xi)=2\int_0^\infty\e^{-ax}\cos(2\pi\xi x)\dd x=\frac{2a}{a^2+4\pi^2\xi^2}$。当$x\ge0$，$\int_\R\e^{-a|t|}\e^{-a|x-t|}\dd t=\frac{\e^{-ax}}{2a}+x\e^{-ax}+\frac{\e^{-ax}}{2a}=\e^{-ax}(x+\frac{1}{a})$；偶延拓得$(f*f)(x)=\e^{-a|x|}(|x|+\frac{1}{a})$。其变换为$[\frac{2a}{a^2+4\pi^2\xi^2}]^2$，可作第二种校验。\par\textbf{方法点评} 连接Fourier余弦变换与卷积性质。\par\textbf{变式与易错点} 本约定下$\widehat{f*g}=\widehat f\,\widehat g$；对Schwartz函数等满足交换条件的函数，还有$\widehat{fg}=\widehat f*\widehat g$，两式均无额外$2\pi$。角频率$\omega=2\pi\xi$的另一约定会改变逆变换及乘积公式的系数。\end{example}

\begin{example}{考研风格｜Laplace解受迫初值问题}{难度：$\star\star\star$}\label{deep:ex:integral_complex_deep:laplaceode}\textbf{题目} 解$y''+y=t,\ y(0)=1,y'(0)=0$。\par\textbf{思路} 导数变换包含初值，拆成标准逆变换。\par\textbf{详细解答} 令$Y=\mathcal Ly$，则$(s^2Y-s)+Y=\frac{1}{s^2}$，故$Y=\frac{s}{s^2+1}+\frac{1}{s^2(s^2+1)}=\frac{s}{s^2+1}+\frac{1}{s^2}-\frac{1}{s^2+1}$。逆变换得$y=\cos t+t-\sin t$。第二法待定系数取特解$t$，齐次解$A\cos t+B\sin t$，初值得$A=1,B=-1$。\par\textbf{方法点评} 本题待定系数更短；Laplace对分段外力和积分方程更统一。\par\textbf{变式与易错点} 漏写$-sy(0)-y'(0)$会得到错误齐次部分。\end{example}
